\chapter{Machine Learning para la predicción de \texorpdfstring{$\boldsymbol{\mu}$}{μ}}
\label{ch:ml}

El Capítulo~\ref{ch:limitaciones} demostró que la estimación de
$\boldsymbol{\mu}$ es el punto más débil del modelo de Markowitz: la media
muestral tiene un error estándar comparable al propio parámetro, y ningún
estimador de $\boldsymbol{\Sigma}$ puede compensar una mala estimación de las
rentabilidades esperadas. Este capítulo explora si las técnicas de Machine
Learning pueden mejorar la estimación de $\boldsymbol{\mu}$ explotando
relaciones predecibles entre las características observables de los activos
y sus rentabilidades futuras.

%--------------------------------------------------------------------
\section{Planteamiento del problema predictivo}
%--------------------------------------------------------------------

\subsection{¿Por qué Machine Learning?}

La estimación tradicional de $\boldsymbol{\mu}$ como media histórica asume
implícitamente que las rentabilidades son independientes e idénticamente
distribuidas (i.i.d.). Si esta hipótesis fuera cierta, la media muestral sería
el mejor estimador lineal insesgado de $\mu_i$. Sin embargo, existe
abundante evidencia empírica de que las rentabilidades no son i.i.d.:
presentan \textbf{autocorrelación} a corto plazo (momentum, reversión a la
media), dependencia de la \textbf{volatilidad} (clustering) y sensibilidad a
\textbf{factores de riesgo} comunes \cite{Cont2001}.

Si existe una función $f$ que relaciona características observables
$\mathbf{x}_{i,t}$ del activo $i$ en el instante $t$ con su rentabilidad
futura $r_{i,t+1}$,
\begin{equation}
    r_{i,t+1} = f(\mathbf{x}_{i,t}) + \varepsilon_{i,t+1},
    \label{eq:ml-model}
\end{equation}
entonces un modelo de ML puede aproximar $f$ a partir de datos históricos,
produciendo una estimación de $\mu_{i,t+1} = \E[r_{i,t+1} \mid
\mathbf{x}_{i,t}]$ potencialmente más precisa que la media incondicional.

La clave está en que el modelo aprende de \textbf{todos los activos}
simultáneamente (o al menos comparte estructura entre ellos), lo que mitiga
la escasez de datos por activo: en lugar de estimar una media
por activo con 36 observaciones, el modelo estima una función $f$ con
$30 \times 36 = 1080$ observaciones, reduciendo la varianza del estimador.

\subsection{Configuración del problema}
\label{subsec:config-ml}

Para cada activo $i \in \{1, \dots, n\}$ y cada mes $t$, definimos:

\begin{itemize}
    \item \textbf{Target:} $y_{i,t} = r_{i,t+1}$, la rentabilidad logarítmica
          del activo $i$ en el mes $t+1$.
    \item \textbf{Features:} $\mathbf{x}_{i,t} \in \mathbb{R}^p$, un vector
          de $p$ características del activo $i$ calculadas exclusivamente con
          información disponible hasta el mes $t$ (inclusive).
\end{itemize}

El modelo se entrena de forma \textbf{agrupada} (\textit{pooled}): se estima una
única función $f$ común a todos los activos, apilando las observaciones de los
$n$ activos de la ventana. Así cada ajuste dispone de $n \times 36$
observaciones en lugar de 36, lo que reduce drásticamente la varianza del
estimador; las predicciones siguen siendo individualizadas porque cada activo
entra con sus propias features $\mathbf{x}_{i,t}$.

%--------------------------------------------------------------------
\section{Ingeniería de características (features)}
%--------------------------------------------------------------------

La selección de features está guiada por la literatura sobre predictibilidad
de rentabilidades \cite{Gu2020} y por la restricción fundamental de este
trabajo: solo se utilizan datos de precio (no fundamentales), ya que el
objetivo es un modelo puramente técnico que pueda automatizarse sin depender
de fuentes de datos externas.

\subsection{Features basadas en rentabilidades pasadas}

\begin{itemize}
    \item \textbf{Rezagos simples:} $r_{i,t}$, $r_{i,t-1}$, $r_{i,t-2}$.
          Capturan efectos de autocorrelación a muy corto plazo (reversión
          o continuación).
    \item \textbf{Momentum:} rentabilidad acumulada en los últimos 3, 6 y 12
          meses:
          \begin{equation}
              \text{mom}_{k}(i,t) = \sum_{j=0}^{k-1} r_{i,\,t-j},
              \quad k \in \{3, 6, 12\}.
              \label{eq:momentum}
          \end{equation}
          El momentum a 12 meses es uno de los factores más robustos en la
          literatura de \textit{cross-section} de rentabilidades
          \cite{Jegadeesh1993}.
    \item \textbf{Volatilidad realizada:} desviación típica de las
          rentabilidades de los últimos 12 meses, $\sigma_{i,t}^{(12)}$.
    \item \textbf{Beta al mercado:} pendiente de la regresión de $r_i$ sobre
          $r_m$ (SPY) en ventana de 36 meses,
          $\beta_{i,t} = \Cov_{36\text{m}}(r_i, r_m)/\Var_{36\text{m}}(r_m)$.
\end{itemize}

\subsection{Features del mercado}

Además de las características propias de cada activo, se incluyen features
que capturan el contexto de mercado:
\begin{itemize}
    \item \textbf{Rezagos del mercado:} $r_{m,t}$, $r_{m,t-1}$ (rentabilidad
          del SPY).
    \item \textbf{Volatilidad del mercado:} $\sigma_{m,t}^{(12)}$, desviación
          típica del SPY en 12 meses.
\end{itemize}

En total, el vector de features tiene dimensión $p = 11$. La
Tabla~\ref{tab:features} resume las features utilizadas.

\begin{table}[htbp]
    \centering
    \caption{Features utilizadas para la predicción de $r_{i,t+1}$.}
    \label{tab:features}
    \begin{tabular}{ll}
        \toprule
        \textbf{Feature} & \textbf{Descripción} \\
        \midrule
        \texttt{r\_lag1}, \texttt{r\_lag2}, \texttt{r\_lag3} & Rezagos del activo ($t$, $t-1$, $t-2$) \\
        \texttt{mom3}, \texttt{mom6}, \texttt{mom12} & Momentum a 3, 6 y 12 meses \\
        \texttt{vol12} & Volatilidad realizada a 12 meses \\
        \texttt{beta\_36m} & Beta rolling al SPY (36 meses) \\
        \texttt{r\_mkt\_lag1}, \texttt{r\_mkt\_lag2} & Rezagos del SPY \\
        \texttt{r\_mkt\_vol12} & Volatilidad del SPY a 12 meses \\
        \bottomrule
    \end{tabular}
\end{table}

%--------------------------------------------------------------------
\section{Modelos de Machine Learning}
%--------------------------------------------------------------------

Se implementan tres modelos que representan distintos compromisos entre sesgo
y varianza, flexibilidad e interpretabilidad.

\subsection{Ridge (regularización L2)}

Ridge \cite{HoerlKennard1970} es una regresión lineal con penalización en la
norma $\ell_2$ de los coeficientes:
\begin{equation}
    \hat{\boldsymbol{\beta}}_{\text{Ridge}} =
    \arg\min_{\boldsymbol{\beta}} \left\{
    \sum_{i,t} (y_{i,t} - \mathbf{x}_{i,t}^T \boldsymbol{\beta})^2
    + \alpha \sum_{j=1}^{p} \beta_j^2
    \right\},
    \label{eq:ridge}
\end{equation}
donde la primera suma recorre las observaciones apiladas de todos los activos
(la muestra agrupada de la Sección~\ref{subsec:config-ml}).

La penalización $\ell_2$ contrae los coeficientes hacia cero sin anularlos:
reduce la varianza del estimador a cambio de un pequeño sesgo, y es adecuada
cuando muchas features tienen poder predictivo débil pero no nulo o están muy
correlacionadas entre sí. En ese caso, una regresión sin penalizar daría
coeficientes muy grandes e inestables (la matriz $\mathbf{X}^T\mathbf{X}$ queda
mal condicionada, casi singular); la penalización de Ridge los estabiliza. El
hiperparámetro $\alpha$ ---de $0$ (mínimos cuadrados ordinarios) a $\infty$
(todos los coeficientes nulos)--- se selecciona por validación cruzada temporal
(\textit{TimeSeriesSplit}) en 10 valores entre $10^{-3}$ y $10^{2}$. Como el
entrenamiento es agrupado sobre activos de escalas heterogéneas, las features se
estandarizan (media 0, varianza 1) usando solo la muestra de entrenamiento.

\subsection{Lasso (regularización L1)}

Lasso \cite{Tibshirani1996} sustituye la penalización $\ell_2$ por $\ell_1$:
\begin{equation}
    \hat{\boldsymbol{\beta}}_{\text{Lasso}} =
    \arg\min_{\boldsymbol{\beta}} \left\{
    \sum_{i,t} (y_{i,t} - \mathbf{x}_{i,t}^T \boldsymbol{\beta})^2
    + \alpha \sum_{j=1}^{p} |\beta_j|
    \right\}.
    \label{eq:lasso}
\end{equation}

La norma $\ell_1$ produce soluciones \textbf{sparse} (muchos coeficientes
exactamente cero): Lasso hace selección automática de features y es más
interpretable que Ridge, a costa de más inestabilidad con features muy
correlacionadas \cite{Zou2005}. Como en Ridge, las features se estandarizan y
$\alpha$ se elige por validación cruzada temporal en el mismo rango.

\subsection{Random Forest}

Random Forest es un método no paramétrico basado en un \textit{ensemble} de
árboles de decisión \cite{Breiman2001}. Cada árbol particiona recursivamente
el espacio de features para minimizar el error cuadrático medio en cada nodo,
y el bosque promedia las predicciones de todos los árboles.

Sus principales ventajas aquí son que captura relaciones no lineales e
interacciones entre features sin especificarlas, es robusto a outliers y a
features irrelevantes, y el promediado sobre muchos árboles (100 en la
implementación) reduce la varianza sin aumentar el sesgo.

Para evitar el sobreajuste se limitan la profundidad máxima de cada árbol a 3
y el número mínimo de muestras por hoja a 5. Estos hiperparámetros, junto con
el número de árboles (100), se mantienen fijos para todo el backtesting. A
diferencia de los modelos lineales, Random Forest no requiere estandarizar las
features, al ser invariante a transformaciones monótonas de cada variable.

%--------------------------------------------------------------------
\section{Validación temporal (walk-forward)}
%--------------------------------------------------------------------

\subsection{El peligro del look-ahead}

En la predicción de series temporales financieras, la validación cruzada
estándar (aleatoria) es metodológicamente incorrecta: al barajar las
observaciones, se rompe la estructura temporal y el modelo podría entrenarse
con datos del futuro para predecir el pasado. Esto produce estimaciones
irrealmente optimistas del rendimiento predictivo y, lo que es más grave,
contamina la selección de hiperparámetros.

Por esta razón, toda la validación se realiza en régimen walk-forward:
en cada mes $t$, el modelo solo puede acceder a datos hasta $t$ para
entrenar, validar hiperparámetros y predecir.

\subsection{Protocolo de backtesting}

En cada mes $t$, el modelo se entrena solo con datos hasta $t$ (ventana de 36
meses), predice $\hat{\mu}_{i,t+1}$ para cada activo a partir de sus features de
$t$, y esas predicciones ---junto con $\hat{\boldsymbol{\Sigma}}$--- alimentan la
optimización de la cartera de máximo Sharpe, cuya rentabilidad se evalúa en
$t+1$. El ciclo se repite 143 veces (enero 2013 -- diciembre 2024). El
Capítulo~\ref{ch:metodologia} detalla el protocolo completo y las medidas que
se toman para evitar el look-ahead.

\subsection{Métricas de evaluación predictiva}

Para evaluar la calidad de las predicciones de $\boldsymbol{\mu}$, se calculan
las siguientes métricas sobre las predicciones out-of-sample de cada modelo:

\begin{itemize}
    \item \textbf{RMSE} (Root Mean Square Error):
          $\sqrt{\frac{1}{nT} \sum_{i,t} (r_{i,t+1} - \hat{\mu}_{i,t})^2}$.
    \item \textbf{MAE} (Mean Absolute Error):
          $\frac{1}{nT} \sum_{i,t} |r_{i,t+1} - \hat{\mu}_{i,t}|$.
    \item \textbf{$R^2$ out-of-sample}:
          $1 - \frac{\sum_{i,t} (r_{i,t+1} - \hat{\mu}_{i,t})^2}
          {\sum_{i,t} (r_{i,t+1} - \bar{r})^2}$,
          donde $\bar{r}$ es la media histórica. Un $R^2$ positivo indica que
          el modelo ML mejora sobre la media histórica.
    \item \textbf{Correlación de Pearson} entre $\hat{\boldsymbol{\mu}}$ y
          $\mathbf{r}_{t+1}$ realizada, agregada a nivel de mes.
\end{itemize}

En finanzas, incluso un $R^2$ pequeño
($\approx 0{,}01$--$0{,}05$) puede ser económicamente significativo si se
traduce en una mejora del ratio de Sharpe de la cartera optimizada.

%--------------------------------------------------------------------
\section{Conexión con la optimización de carteras}
%--------------------------------------------------------------------

El objetivo no es minimizar el RMSE de
$\hat{\boldsymbol{\mu}}$, sino maximizar el rendimiento out-of-sample de la
cartera: un RMSE bajo no garantiza una buena cartera, y el QP puede
beneficiarse más de acertar el \textit{orden} (ranking) de los activos que de
minimizar el error cuadrático. Por eso la evaluación definitiva se hace sobre
las métricas de cartera (Sharpe, drawdown, turnover) del
Capítulo~\ref{ch:resultados}, y las predictivas se reportan solo como
diagnóstico complementario.

%--------------------------------------------------------------------
\section{Síntesis}

En resumen, este capítulo aborda la principal limitación del modelo clásico
---la estimación de $\boldsymbol{\mu}$--- con un marco de Machine Learning que
usa 11 features exclusivamente de precio, tres modelos complementarios (Ridge,
Lasso y Random Forest), una validación temporal estricta sin look-ahead y, como
métrica final, el rendimiento de la cartera y no el RMSE de la predicción. El
Capítulo~\ref{ch:metodologia} detalla cómo se integran estos modelos con los
estimadores de covarianza del Capítulo~\ref{ch:limitaciones} en el backtesting
walk-forward.
